\subsection*{从一个已处理标记推导完成记录}

先考虑只保存 S 和 \code{handled=true} 的方案。R2 找到标记，知道服务处理过，却无法知道是成功还是余额不足；即使旁边记了 \code{succeeded}，也还找不到 O1 和 D1。若随后重新建一张订单来凑响应，便会增加新的业务记录；若按商品与时间搜索一条看起来相近的订单，又可能把 Alice 的另一笔合法购买认成 P1。完成记录需要直接保存能精确恢复这项操作的关联。

可以把处理过标记遗漏的信息写成两个具体历史。历史 H1 从 10000 开始，S1 先扣 3000 得到 P1、D1 和当次 7000，另一合法 S2 再扣 1000，最后为 6000。历史 H2 也从 10000 开始，S2 先扣 1000 得到当次 9000，S1 再扣 3000，最后同样为 6000。两段历史中，S1 的商品、数量、Q7 与金额完全相同，\code{S1.handled} 都为 true，账户最终都为 6000；应恢复给 S1 的当次余额却分别为 7000 与 6000。

\begin{figure}[htbp]
\centering
\begin{tikzpicture}[x=1cm,y=1cm]
\node[booknote] at (0,.6) {同一初始余额 10000};
\node[booknode,text width=4.2cm] (a) at (0,-.4) {H1：S1 扣 3000 得 7000\\随后 S2 扣 1000 得 6000};
\node[booknode,text width=4.2cm] (b) at (0,-2.3) {H2：S2 扣 1000 得 9000\\随后 S1 扣 3000 得 6000};
\node[booknode,text width=3cm] (o) at (5,-1.3) {共同可读资料\\S1 已处理、原输入\\当前余额 6000};
\node[bookresult] (r1) at (9,-.4) {S1 应返回 7000};
\node[bookresult] (r2) at (9,-2.3) {S1 应返回 6000};
\draw[bookarrow] (a)--(o);
\draw[bookarrow] (b)--(o);
\draw[bookalert] (o)--(r1);
\draw[bookalert] (o)--(r2);
\end{tikzpicture}
\caption{只保存处理标记与当前余额，会把需要不同原结果的两条历史混为一谈。}
\label{fig:ch05-indistinguishable}
\end{figure}


若恢复算法只读取处理过标记、S1 输入与当前余额，它在两段历史中取得完全相同的数据，便只能作出同样的回答，至少有一段历史的原返回值会答错。这是观察资料不足导致的不可区分：需要回答的问题区分两段历史，算法能读取的资料没有区分它们。保存 D1 当次余额以及到 D1 的确定关联，便为两段历史留下不同证据。若流水本来已经保存这些内容，完成记录可以引用它，但恢复必须沿该引用读取，不能仍只使用标记与当前余额。

为了明确这个限制，把恢复者能读取的信息记为 \(O(H)\)，其中 H 表示一段执行历史；把它在历史 H 上应为 S1 恢复的原业务结果记为 \(R_H(S_1)\)。只依赖这些信息的恢复算法记为 \(g\)。上例给出了
\[
 O(H_1)=O(H_2),\qquad R_{H_1}(S_1)\ne R_{H_2}(S_1).
\]
由于 \(g\) 接到的资料相同，\(g(O(H_1))\) 与 \(g(O(H_2))\) 也相同，便无法同时等于两份不同的原结果。这里的差别具体落在 7000 与 6000，未要求恢复者猜出不可观察的事件先后。要消除这种困难，需要增加足以区分结果的持久资料，或收窄接口本来就不提供的历史字段。

连是否成功也有类似反例。一段历史中 S1 成功扣 3000、随后另一笔充值 3000；另一段历史中 S1 因业务规则拒绝、没有后续变动。若两个历史起始余额相同，最终余额也相同，\code{handled=true} 也相同。只有终态或可确定关联的业务事实能把它们分开。这个反例并不要求系统故意遗失已有流水，它说明一旦设计宣称只靠处理标记恢复，必须检查标记实际承载了多少结果信息。

完成记录在本章指一条已提交操作所留下的身份、输入绑定、终态和原业务结果。它可以全部位于 \code{purchases} 一行，也可以通过不可变引用指向同一事务中保存的结果。第四章已经把 S 放在 \code{purchases}，本章继续采用这一路径，并补全它保存的内容。购买记录既承担受理身份，又描述业务处理；成功结果引用订单和流水，拒绝结果也仍具有自己的 \code{purchase\_id}。

\begin{table}[htbp]
\centering\small
\begin{tabularx}{\linewidth}{@{}p{2cm}p{4.8cm}X@{}}
\toprule
字段组 & P1 示例 & 恢复时的用途\\
\midrule
完整身份 & 账户 1、动作 v1、E1、K1 & 找到同一操作并限制重复受理\\
输入绑定 & purchase-input-v1 与原规范化字段 & 按原规则判断同参或异参\\
业务终态 & succeeded & 说明处理已作出最后结论\\
业务关联 & P1、O1、D1 & 沿确定引用找到购买、订单与流水\\
原成交 & paper-001、数量 1、Q7、CNY 3000 & 解释当时接受的商品与金额\\
原结果表示 & purchase-result-v1 与结果字段 & 升级后仍按原含义返回\\
时间边界 & 受理时点、结果记录时点、恢复截止 & 解释处理历史并限制旧请求\\
\bottomrule
\end{tabularx}
\caption{完成记录保存恢复原结论所需的身份、内容与解释依据。}
\end{table}

\code{purchase-input-v1} 标记输入解释规则；规范化字段把省略数量与明确数量 1 等约定为相同的写法统一表示。\code{purchase-result-v1} 则标记原结果字段的含义和编码规则。两者分别解决输入比较与输出解释：输入相同，不代表任何软件版本都可以任意改变结果中金额的单位。若 v2 将分改为元，继续把旧整数 3000 解释为元，会把一笔 30 元购买显示为 3000 元。

购买编号在受理时生成，成功与拒绝都可使用。成功购买的记录填写 \code{order\_id} 和 \code{debit\_id}，引用实际建立的订单与扣款；余额不足的记录保存拒绝原因和当次要求的金额，订单与扣款关联留空。恢复时先读终态：成功分支读取成交关系，拒绝分支读取原拒绝原因。两种记录分别具有完整的恢复内容，空关联正是余额不足分支没有成交的表示。

\begin{figure}[htbp]
\centering
\begin{tikzpicture}[x=1cm,y=1cm]
\node[booknode,text width=10cm] (s) at (0,0) {完整 S：账户 1、purchase.create.v1、E1、K1\\绑定原规范化输入与规范化版本};
\node[bookresult] (p) at (0,-1.5) {P1 的唯一终态及结果版本};
\node[booknode,text width=4.2cm] (a) at (-3,-3.2) {成功\\O1 / D1 / CNY 3000\\当次余额 7000};
\node[booknode,text width=4.2cm] (b) at (3,-3.2) {余额不足拒绝\\原因及所需 CNY 3000\\无订单、无扣款};
\draw[bookarrow] (s)--(p);
\draw[bookarrow] (p)--(a);
\draw[bookarrow] (p)--(b);
\end{tikzpicture}
\caption{完成记录把操作身份、原输入和该终态所需的完整结果连接起来。}
\label{fig:ch05-completion-record}
\end{figure}


记录时间还要避免过度命名。服务可以在准备终态时写 \code{result\_recorded\_at}，表示这份结果在事务内被记录的时间；它随后才发送 COMMIT。若字段叫 \code{committed\_at}，读者容易把服务取到的时钟值当成数据库实际提交时刻。简单实现可以只保存受理与结果记录时间，不提供精确提交时间。某项业务确实需要数据库提交顺序或时刻时，应选择能够提供这种证据的数据库机制，并说明其精度与来源。

\subsection*{用关系限制编号和金额的错配}

S 的唯一性限制同一身份只有一行购买。订单再以 \code{purchase\_id} 引用购买，并对该列设唯一约束，可限制本例一笔购买最多一张订单。流水也引用 \code{purchase\_id} 和 \code{order\_id}，并按本例的一次全额本地扣款规则，对 \code{purchase\_id} 施加唯一约束，防止 D1 指向不存在的 P1 或 O1。

仅有两条独立外键，仍可能形成 \code{D1.purchase\_id=P1}、\code{D1.order\_id=O2}，而 O2 属于 P2。两个被引用对象都存在，组合关系却错了。可以在订单中为 \code{(order\_id,purchase\_id)} 建立可被引用的唯一组合，并使流水的这两个字段共同引用它；也可以通过受控写入与检查维持同样关系。这样检查 D1 时，数据库要求它引用的 O1 确实属于 P1。本例用共同事务构造正确引用，数据约束再减少遗漏路径。

第四章已经从订单回填并加上流水的 \code{purchase\_id}。现在继续执行下面的约束演进，保留第一章已有的订单外键与订单唯一关系：

\begin{BookCode}
ALTER TABLE orders ADD CONSTRAINT order_purchase_pair
  UNIQUE (order_id, purchase_id);
ALTER TABLE debits
  ADD CONSTRAINT debit_one_per_purchase UNIQUE (purchase_id),
  ADD CONSTRAINT debit_order_purchase_pair
    FOREIGN KEY (order_id, purchase_id)
    REFERENCES orders (order_id, purchase_id);
\end{BookCode}

结果版本与 \code{result\_payload} 两列也已在第四章加入；本章补全它们的写入与读取含义，结构仍是同一组表。成功结果中的订单号和流水号取自本事务的实际写入，拒绝结果不携带这两个编号。

\begin{table}[htbp]
\centering\small
\begin{tabularx}{\linewidth}{@{}>{\raggedright\arraybackslash}p{2.1cm}>{\raggedright\arraybackslash}p{6.1cm}>{\raggedright\arraybackslash}X@{}}
\toprule
表 & 核心列 & 已建立的约束\\
\midrule
\code{accounts} & \code{account\_id}；\code{balance\_minor} & 账户主键；余额非空且非负\\
\code{purchases} & \code{purchase\_id}；完整 S 的四列；\code{normalization\_version}、\code{canonical\_input}；商品、数量、报价、金额、币种；\code{status}、\code{result\_version}、\code{result\_payload}、\code{created\_at} & 购买主键；账户外键；S 四列非空且组合唯一；金额为正、CNY、状态取值检查；终态快照与结果完整\\
\code{orders} & \code{order\_id}、\code{purchase\_id}；商品、金额、币种、\code{status} & 订单主键；购买外键且一购一单；订单/购买组合唯一；金额为正、CNY、paid\\
\code{debits} & \code{debit\_id}、\code{order\_id}、\code{purchase\_id}、\code{account\_id}；金额、币种、\code{balance\_after\_minor}、\code{created\_at} & 流水主键；账户/订单/购买外键；一单一款、一购一款；订单与购买组合外键；金额为正、CNY、当次余额非负\\
\bottomrule
\end{tabularx}
\caption{第一部分四张表的核心列与约束总览}
\label{tab:ch05-final-schema}
\end{table}

表\ref{tab:ch05-final-schema}汇总第一章建表、第四章身份与结果列演进以及本节组合外键后的结构。P1 的订单号、扣款号保存在原结果中，也能沿订单、流水的购买外键查出；购买表没有再增加独立的这两个编号列。成功时三表金额相等、终态后绑定不可变以及内部 \code{pending} 不提交，仍由受控写入协议维持。

\begin{figure}[htbp]
\centering
\begin{tikzpicture}[x=1cm,y=1cm]
\node[booknode] (p1) at (0,0) {购买 P1};
\node[booknode] (p2) at (6,0) {购买 P2};
\node[booknode] (o1) at (0,-1.6) {订单 O1 属于 P1};
\node[booknode] (o2) at (6,-1.6) {订单 O2 属于 P2};
\draw[bookarrow] (p1)--(o1);
\draw[bookarrow] (p2)--(o2);
\node[bookresult,text width=4.2cm] at (0,-3.3) {组合引用 (O1,P1)\\匹配同一行，允许};
\node[booknode,text width=4.2cm] at (6,-3.3) {组合引用 (O2,P1)\\没有匹配行，拒绝};
\draw[bookalert,dashed] (p1)--node[above,font=\footnotesize,sloped] {分别存在仍不足够} (o2);
\end{tikzpicture}
\caption{组合外键检查订单归属，避免把另一笔购买的订单接到 P1。}
\label{fig:ch05-composite-foreign-key}
\end{figure}


金额也要核对：P1 的成交为 CNY 3000，O1 应保留同一成交金额，D1 的支出应为同币种 3000，原结果也返回 3000。外键只检查引用，不会自动证明这些数值相等。单行 CHECK 可以检查金额为正、数量为 1、状态与本行空字段的关系；CHECK 是数据库在写入时判断本行条件的约束。本例需由业务代码在共同事务中使用同一份经验证的 Q7 结果写各表，并限制任意直接改账。

数据重复保存可以帮助恢复，也带来一致性责任。若 \code{purchases.result\_payload} 里写了 D9，实际流水却是 D1，服务即使每次稳定返回 D9，也稳定返回了错误内容。可在提交前用实际插入返回的编号构造原结果，不再从另一套临时变量猜测。对成交金额与当次余额，可以选择以不可变流水字段为结果来源，或在同一事务中把它们复制到结果内容；选择复制时应核对两份数值。

\begin{figure}[htbp]
\centering
\begin{tikzpicture}[x=1cm,y=1cm]
\node[booknode,text width=2.6cm] (p) at (0,0) {P1\\S / Q7 / CNY 3000};
\node[booknode,text width=2.6cm] (o) at (4,0) {O1\\购买编号 P1};
\node[booknode,text width=2.8cm] (d) at (8,0) {D1\\P1 + O1 / 扣 3000};
\draw[bookarrow] (p)--(o);
\draw[bookarrow] (o)--(d);
\node[bookresult,text width=10cm] (r) at (4,-2) {完成结果：P1 / O1 / D1 / Q7 / CNY 3000 / 当次余额 7000};
\draw[bookarrow] (p)--(r);
\draw[bookarrow] (d)--(r);
\end{tikzpicture}
\caption{结果从实际形成的购买、订单和扣款关系中构造。}
\label{fig:ch05-result-relations}
\end{figure}


本章保存一份明确的 \code{result\_payload}，表示按结果版本组织的业务内容，方便恢复时直接取回。它按 JSON 对象保存，金额仍为整数；对象的字段顺序和空白可在序列化时变化。下一节会据此区分业务结果重放与逐字节回放。

终态后，S、输入版本、规范化输入、结果版本和 \code{result\_payload} 都应保持不变。订单现状在 \code{orders} 中更新，账户余额在 \code{accounts} 中更新。可以让普通应用角色只能执行规定的写入路径，限制终态更新，并对维护操作留审计记录；需要强制保护时，再由数据库触发器检查旧行与新行。触发器是数据库在指定写入事件发生时自动执行的检查程序，可拒绝终态后改变身份或原结果。唯一约束本身只限制键重复，不能阻止有人把 K1 改成 K9。

保存正确结果与保持结果不变是两项要求。第一项要核对结果声称的编号、金额、终态是否与业务事实相符；第二项要核对后来的恢复是否继续返回同一份结果。把 D9 错存一次并永久返回，满足不变性却违反真实性；首次正确返回 D1，后来覆盖为 D9，又破坏了不变性。结果字段与实际写入共用来源，加上终态后的受控更新，分别承担这两项责任。

\begin{figure}[htbp]
\centering
\begin{tikzpicture}[x=1cm,y=1cm]
\node[booknote] at (1,1.05) {实际扣款均为 D1};
\node[booknode] at (-3,0) {首次 D1\\恢复 D1};
\node[bookresult] at (-3,-1.5) {真实，且保持不变};
\node[booknode] at (1,0) {首次 D9\\恢复 D9};
\node[booknode] at (1,-1.5) {保持不变，但错误};
\node[booknode] at (5,0) {首次 D1\\恢复 D9};
\node[booknode] at (5,-1.5) {恢复结果发生变化};
\end{tikzpicture}
\caption{结果真实性与结果不变性需要分别检查。}
\label{fig:ch05-truth-stability}
\end{figure}


\subsection*{两个提交之间分别会失去什么}

把扣款先提交、结果稍后保存，看起来可以缩短购买事务，却制造了明确的空隙。事务 A 扣 3000、建立订单和流水，提交后服务退出；事务 B 尚未写 S 到 P1 的映射。R2 查询 S 得到空，再走新购买，余额可能由 7000 扣到 4000。即使后续人工能从 D1 找回旧记录，自动恢复入口当时已经缺少必须遵循的身份关系。

反过来，先提交 S 与 \code{succeeded}，再扣余额，也有对应空隙。R2 读到成功，显示 O1 或付款完成；实际扣款事务尚未成功，账户仍为 10000，订单甚至还不存在。先提交 \code{pending} 则会形成长期已占用但未完成记录，需要额外的任务与接管协议才能继续。本章没有引入这类可见中间状态，\code{pending} 始终只是同一事务内的临时填写状态，任何路径都只提交完整终态或回滚。

RIFL 要求完成记录与操作修改共同可靠保存，并在结果持久保存以后才向调用者返回。本章用单数据库事务具体实现这项关系\cite{rifl}。

\begin{figure}[htbp]
\centering
\begin{tikzpicture}[x=1cm,y=1cm]
\node[booknode,text width=3cm] (a) at (0,0) {先提交扣款\\余额已为 7000};
\node[booknode,text width=3cm] (b) at (6,0) {崩溃：结果未保存\\恢复时缺少结论};
\node[booknode,text width=3cm] (c) at (0,-1.7) {先提交成功结果\\声明 D1 已扣款};
\node[booknode,text width=3cm] (d) at (6,-1.7) {崩溃：扣款未发生\\余额仍为 10000};
\node[bookresult,text width=9cm] at (3,-3.3) {共同事务：身份、扣款、订单、流水\\原结果与终态一起提交};
\draw[bookalert] (a)--node[above,font=\footnotesize] {两次提交之间} (b);
\draw[bookalert] (c)--node[above,font=\footnotesize] {两次提交之间} (d);
\end{tikzpicture}
\caption{分开提交会留下业务事实与完成结果不相符的崩溃窗口。}
\label{fig:ch05-split-commits}
\end{figure}


可以据此给已提交状态写出一项不变量：一份可对外恢复的成功完成记录，必须与它声明的购买、订单、扣款及资金变化共同成立；本例每一份成功业务效果，也必须有能够沿原 S 找到的完成记录。拒绝终态则需要完整身份、输入与拒绝结果共同成立，且没有该购买的订单或扣款。这里的不变量是对允许的已提交状态提出的持续条件，事务内部的 \code{pending} 不在外部观察集合中。

先扣款后记结果的崩溃历史违反第二个方向，留下效果却缺少恢复依据；先记成功后扣款的历史违反第一个方向，留下结果却缺少它声称的效果。共同事务消除了这两类中间状态成为已提交事实的机会，但仍要求事务写入的编号和金额本身正确。它不会替错误代码把 4000 自动纠正为 Q7 的 3000。把原子关系与数值关系分开核对，才能看清事务提供了什么保证、应用还承担什么责任。

\begin{figure}[htbp]
\centering
\begin{tikzpicture}[x=1cm,y=1cm]
\draw[draw=Ink!50,dashed] (-2,-3.3) rectangle (9,1.25);
\node[booknote] at (3.5,.95) {只观察已提交状态};
\node[bookresult,text width=3.4cm] (a) at (0,-.5) {P1 成功完成记录\\声明 O1、D1\\金额 3000};
\node[booknode,text width=3.8cm] (b) at (6,-.5) {对应订单、扣款与历史金额\\共同成立并关联同一 P1};
\draw[bookarrow,<->] (a)--(b);
\node[booknode,text width=8.5cm] at (3.5,-2.5) {P1 终态拒绝：保存原因与原输入；该购买没有订单和扣款};
\end{tikzpicture}
\caption{成功完成记录与其声明的业务事实双向对应；拒绝分支保留自身的完整结论。}
\label{fig:ch05-atomic-invariant}
\end{figure}


同一事务中的执行顺序可以保持为：插入 S 与内部 \code{pending}；核验首次成交依据；条件更新余额；插入 O1 与 D1；从这些写入的实际结果构造 \code{result\_payload}；将 P1 改为 \code{succeeded}；提交。余额不足则填写拒绝终态与对应结果后提交。服务收到明确提交成功以后，才把终态发送出去。中途构造结果失败，也要回滚整笔事务，不能留下扣款再打算下一次补结果。

共同事务提供原子性，即这组数据库变化共同提交或共同撤销。持久性则说明已确认提交的数据在规定故障后仍可恢复，依赖数据库日志、存储与配置。PostgreSQL 18 关于异步提交的官方说明指出，关闭同步提交等待可能允许已经向客户端报告成功的近期事务在崩溃后丢失\cite{pg}。本章的恢复推导假定数据库兑现已确认提交的持久性，也假定未来切换的权威数据库保留这些事实；多个副本和旧备份恢复需要另行分析。

\subsection*{COMMIT 发出以后断开连接}

现在从初始 10000、空记录开始，A 已在事务内完成 P1、O1、D1 和原结果，发送 COMMIT 后连接中断。服务没有收到数据库的确认。此时服务所知道的是提交结果未知：数据库可能已经提交，也可能没有提交，还可能正在结束事务。连接异常只描述服务观察到的通信故障，不给出这三者的选择。

如果服务把异常写成一个已保存的失败结果，还需回答它通过哪个事务写入。另开事务覆盖 S，可能与原成功事务冲突；不核对原状态就写 \code{failed}，会把已经成立的成功改成失败。若直接返回余额不足，则更没有业务依据。合理的暂时响应应说明结果尚未确认，要求继续原身份恢复，同时保留本次 R1 的诊断记录。

服务可以先自己恢复，客户端也可以随后发 R2。恢复者丢弃失效连接，在新的连接中使用原 S、原输入和同一套受约束受理流程。原事务已经提交时，唯一插入返回零行，下一条 SELECT 读到 P1；原事务回滚时，插入可以成功，新事务才执行购买；原事务还在结束时，唯一检查等待。候选购买号 P2 可以在恢复中生成，但只有实际插入成功时才可用于业务，读取原记录的分支返回 P1。

\begin{figure}[htbp]
\centering
\begin{tikzpicture}[x=1cm,y=1cm]
\node[booknode,text width=8cm] (a) at (0,0) {COMMIT 已发出，确认丢失；应用暂时不知道数据库结局};
\node[bookresult] (b) at (0,-1.3) {新连接携带原 S 进入唯一插入};
\node[booknode,text width=3cm] (c) at (-4,-3) {原事务已提交\\插入零行\\下一 SELECT 读 P1};
\node[booknode,text width=3cm] (d) at (0,-3) {原事务已回滚\\插入成功\\才进入购买分支};
\node[booknode,text width=3cm] (e) at (4,-3) {原事务尚未结束\\唯一检查等待\\或暂时错误后恢复};
\draw[bookarrow] (a)--(b);
\draw[bookarrow] (b)--(c);
\draw[bookarrow] (b)--(d);
\draw[bookarrow] (b)--(e);
\end{tikzpicture}
\caption{提交确认丢失后，恢复仍由原身份的唯一约束与已提交记录决定。}
\label{fig:ch05-unknown-commit}
\end{figure}


这里遵循第 4.3 节的两条语句规则：插入返回零行以后，下一条独立 SELECT 按完整 S 读取原记录，不把冲突后的读取压入同一条 WITH 语句\cite{pg}。

若只想先查询，找到 P1 足以恢复结果；查不到时仍需原身份的受约束重试。它不能据此认定原事务回滚。恢复读取还要求走保存该身份的权威数据路径，按完整 S 查询，没有把已完成但已退款等记录过滤掉，没有提前清理记录。异步副本可能暂时缺少 P1，按订单未关闭之类当前状态过滤也可能漏掉 P1；这些空结果都不能打开不带身份的新扣款路径。

\begin{figure}[htbp]
\centering
\begin{tikzpicture}[x=1cm,y=1cm]
\node[booknode] (s) at (0,0) {完整 S};
\node[booknode] (a) at (3,0) {正确权威数据};
\node[booknode] (r) at (6.5,0) {仍保留的完成记录};
\node[bookresult] (v) at (10,0) {正确版本解释};
\draw[bookarrow] (s)--(a);
\draw[bookarrow] (a)--(r);
\draw[bookarrow] (r)--(v);
\node[booknode,text width=3cm] (f) at (3,-2) {旧副本读空};
\node[booknode,text width=3cm] (g) at (6.5,-2) {按现状过滤\\或提前清理};
\draw[bookalert,dashed] (a)--(f);
\draw[bookalert,dashed] (r)--(g);
\node[booknote,text width=10cm] at (5,-3.1) {这些假空都不能证明原购买没有提交，因而不能据此重新扣款};
\end{tikzpicture}
\caption{记录存在与恢复路径能够找到并解释记录，是两个需要同时满足的条件。}
\label{fig:ch05-evidence-path}
\end{figure}


完成记录的保存和可找到性在这里共同起作用。数据库中某处存在 P1，只能说明资料还在；恢复请求必须被路由到拥有该事实的数据范围，使用相同身份列，并且在承诺期限内仍能读取其版本。若把原结果另存到归档库，却让重试入口只查空的在线表，持久性本身仍然成立，恢复协议已经失去依据。归档、分片或副本设计需要把查找规则一起迁移；在尚未确认旧事实之前，入口继续按原身份仲裁或报告未知。

如果数据库长期不可用，恢复就暂时无法取得终态。服务可以控制每次等待时长和重试次数，返回可恢复的暂时错误，避免一个 HTTP 连接无限占用资源。此时 S 与已确认输入继续保存在客户端或恢复任务里。重复尝试可以增加获知结果的机会，其进展仍取决于数据库恢复、网络可达和结果记录可读，唯一约束不会自行使不可用系统恢复。
\FloatBarrier
